\section{Sharp lower bounds at every source count}
\label{sec:embedding}

The restriction theorem is an upper argument. To obtain a lower witness at
an arbitrarily prescribed source count, extra states must not create cheap
merges. The globally positive curvature floor of the inherited endpoint
construction permits a direct exterior-anchor embedding. Its auxiliary
coordinates may be arbitrarily far apart; an affine change places every
posterior in the required interior interval afterward.

\begin{theorem}[Uniform-in-$N$ matching lower witnesses]
\label{thm:embedding}

There are absolute constants $c>0$ and $r_0$ such that, for every integer
$N\ge4$ and every integer $r\ge r_0$, there is a source with exactly
$N$ distinct posteriors in $[1/4,3/4]$, all source weights positive,
and a polynomial generator with positive curvature of degree at most
$r$, whose induced smooth strictly proper losses are bounded in
$[0,1/2]$ on the full report interval $[0,1]$, for which
\[
 \rho_{\mathrm{det}}^{(N-2,N-1)}
 \ge\rho_{\mathrm{stoch}}^{(N-2,N-1)}
 \ge c\frac{r^2}{(\log r)^2}.
\]
For example, $c=1/9216$ suffices with a suitably large absolute $r_0$.
The masses, generator, and posterior locations may depend on both $N$
and $r$; no minimum posterior gap is asserted.


The construction also preserves the original quartet deterministic price
exactly; no exact stochastic price equality is asserted.
\end{theorem}
\begin{proof}
\par\medskip\noindent\textit{Construction and anchor penalty.} 
Use the endpoint quartet of Appendix~\ref{sec:ext-weights}, writing
\[
 d=(4\log n/n)^2,\quad b=d^2,\quad a=1/2-b,
 \quad (x_A,x_B,x_C,x_D)=(0,2d,1-2d,1),
\]
and $g_n=\varphi_n''=\kappa_n(t)+\kappa_n(1-t)+b$.
In particular $g_n(t)\ge b>0$ for every real $t$, and its degree is
exactly $2n$.
Let $N\ge5$, $m=N-4$, and $\lambda=1/2$. Give the original four
states masses $\lambda(a,b,b,a)$ and give $m$ additional anchors equal
masses $\eta=1/(2m)$. For all sufficiently large $n$, the smallest
of these masses is
\[
 \mu=\min\{\lambda b,\eta\}>0.
\]
Set
\[
 L=\frac2{\sqrt{b\mu}},\qquad
 z_j=-jL\quad(1\le j\le m).
\]
This is an auxiliary real coordinate; the legal posterior coordinates
are defined below. Every pair involving an anchor has gap at least $L$.
The global curvature floor yields the strong-convexity Jensen estimate
\begin{align*}
 D(uv)
 &\ge\frac b2\frac{p_up_v}{p_u+p_v}(u-v)^2\\
 &\ge\frac{b\min(p_u,p_v)}4(u-v)^2
 \ge\frac{b\mu L^2}4=1.
\end{align*}
The first inequality follows by subtracting the convex quadratic
$bx^2/2$ from the generator and applying Jensen to the remaining
convex function. The other inequalities use
$p_up_v/(p_u+p_v)\ge\min(p_u,p_v)/2$.
Thus the anchor penalty holds at every sufficiently large $n$ admitted
by the original quartet, with no additional $N$-dependent threshold.

\par\medskip\noindent\textit{Top-capacity obstruction.} 
Let $P,q,T,J$ denote the original quartet costs. Equations~\eqref{eq:ext-nearpair}--\eqref{eq:ext-coarse} give
\[
 bd/2\le P\le4bd,\qquad0<q\le bd^2,
 \qquad T\ge b/4,\qquad J\ge b/2,
\]
and every core pair other than $BC$ costs at least $P$.
The augmented flat optima satisfy
\[
 0<\operatorname{OPT}_{N-1}\le\lambda q,
 \qquad0<\operatorname{OPT}_{N-2}\le2\lambda P\le8\lambda bd,
\]
by keeping all anchors singleton and using the quartet's fine or coarse
candidate. Refine any deterministic hierarchy to use exactly $N-1$
fine cells and $N-2$ coarse cells; refinement cannot increase costs.
The fine partition merges one pair. If it is not $BC$, its cost is at
least $\lambda P$: this holds for every other core pair, and an
anchor pair costs at least one, exceeding $\lambda P$ for sufficiently
large $n$. Its fine ratio is then at least $P/q\ge1/(2d)$.

If the fine pair is $BC$, the coarse partition merges two fine blocks.
A merge among core blocks has cost at least $\lambda b/4$ by the
three original quartet coarsenings. Any merge involving an anchor
contains a pair of cost at least one, hence also has cost at least one
by the Jensen decomposition $D(S)\ge D(U)$ for $U\subset S$.
Therefore its coarse ratio is at least
\[
 \frac{\lambda b/4}{8\lambda bd}=\frac1{32d}.
\]
This proves the deterministic lower bound $1/(32d)$ for all $N\ge5$;
the $N=4$ statement is the original quartet. Lemma~\ref{lem:ext-rounding} gives the stochastic lower bound $1/(64d)$ for every $N$.
No assertion that optimal cells are contiguous is needed.

\par\medskip\noindent\textit{In fact, the deterministic quartet price is preserved exactly.} 
Let $q_4,Q_4$ be the actual quartet flat optima at capacities three and
two, and let $C$ be the cost of merging the entire scaled core into
one cell. For a cell of total mass $\lambda$ supported on $[0,1]$,
its Jensen kernel satisfies
\[
 K(t)=\min\!\left\{\sum_i p_i(t-x_i)_+,
                       \sum_i p_i(x_i-t)_+\right\}
 \le\lambda\min(t,1-t)\le\lambda/2.
\]
Since $\int_0^1g_n=2+b$ and $\lambda=1/2$, we have
$C\le(2+b)/4<1$ for all sufficiently large $n$.
An augmented flat partition with a nonsingleton cell containing an anchor
costs at least one,
whereas a core-only candidate costs at most $C$. Thus
\[
 \operatorname{OPT}_{N-1}=\lambda q_4,\qquad
 \operatorname{OPT}_{N-2}=\lambda Q_4.
\]
Choose a core fine optimum and any compatible two-cell core coarsening.
Its augmented hierarchy price is at most $C/(\lambda Q_4)$, because
its fine ratio is one and $C\ge\lambda Q_4$. Any hierarchy whose
coarse level mixes an anchor has ratio at least
$1/(\lambda Q_4)>C/(\lambda Q_4)$; an anchor merged at the fine level
is also merged at the coarse level. Therefore every optimal augmented
deterministic hierarchy leaves every anchor singleton at both levels.
Its remaining core hierarchy has capacities two and three, proving
\[
 \rho_{N,\mathrm{det}}=\rho_{4,\mathrm{det}}
\]
for this embedding. This equality is special to the constructed sources;
the upper restriction theorem for an arbitrary source asserts only
domination. No corresponding exact stochastic equality is asserted.

\par\medskip\noindent\textit{Interior support and bounded loss.} 
The auxiliary source hull is $[-mL,1]$. Define
\[
 s=2(1+mL),\qquad t=-mL-s/4,
 \qquad x=sy+t,
 \qquad M=\max_{x\in[t,t+s]}g_n(x).
\]
Then $y=1/4+(x+mL)/[2(1+mL)]$ maps the entire hull onto
$[1/4,3/4]$. The transformed generator
\[
 \widetilde\varphi_n(y)=\frac{\varphi_n(sy+t)}{s^2M}
\]
has curvature $\widetilde g(y)=g_n(sy+t)/M\in(0,1]$ throughout
$[0,1]$, with unchanged degree $2n$. Every augmented-source distortion
is multiplied by the same positive factor $1/(s^2M)$, so all prices
are preserved. Symmetry is not needed: define
\[
 \ell(q,1)=\int_q^1(1-v)\widetilde g(v)\,dv,
 \qquad
 \ell(q,0)=\int_0^q v\widetilde g(v)\,dv.
\]
Both losses lie in $[0,1/2]$ on all of $[0,1]$, are smooth, and the
derivative of conditional risk is $(q-p)\widetilde g(q)$ at posterior
$p$. Therefore the loss is strictly proper and its excess risk is
$B_{\widetilde\varphi_n}(p,q)$.

Finally, put $n=\lfloor r/2\rfloor$. For sufficiently large $r$,
$n\ge r/3$ and $\log n\le\log r$, so
\[
 \frac1{64d_n}
 =\frac{n^2}{1024(\log n)^2}
 \ge\frac{r^2}{9216(\log r)^2}.
\]
All large-degree requirements above come from the original quartet or
absolute inequalities such as $a\ge b$ and $\lambda P<1$, proving
that the same $r_0$ works for every finite $N$. There is no claim of
a uniform positive gap or of nonvanishing absolute excess risk after
normalization.

For explicit absolute-risk bookkeeping, the normalized flat costs obey
\[
 0<\widetilde{\operatorname{OPT}}_{N-1}
 \le\frac{\lambda d^4}{s^2M},\qquad
 0<\widetilde{\operatorname{OPT}}_{N-2}
 \le\frac{8\lambda d^3}{s^2M}.
\]
The interval $[t,t+s]$ contains $[0,1]$, and
$\int_0^1g_n=2+b$, so $M\ge2+b$. Also $s\ge2$ and
$\lambda=1/2$. Thus for $N\ge5$ the respective bounds are at most
$d^4/16$ and $d^3/2$, and both tend to zero uniformly in $N$.
The raw anchor penalty of one is multiplied by the same
$1/(s^2M)$; it is not a positive absolute penalty after normalization.
For $N=4$, the normalization in Appendix~\ref{sec:ext-weights} gives the
respective bounds $d^4/8$ and $d^3$. These weaker constants therefore
give vanishing flat costs uniformly over all $N\ge4$ as well.

\end{proof}
